%% ======================================================================
%% chap6/evaluation_main.tex
%%
%% Chapter 6: Evaluation and Testing
%%
%% Describes the evaluation methodology used to assess CCRMS against the
%% research questions: the four-layer testing architecture, the eight
%% test scenarios, seven automated performance benchmarks, the eleven
%% KPIs, testing infrastructure, and validity considerations. Results
%% themselves are reported in Chapter 7; this chapter covers HOW the
%% system was evaluated, not WHAT was measured.
%%
%% Sections:
%%   6.1 Evaluation Approach (Table 6.1)
%%   6.2 Testing Layers
%%   6.3 Performance Testing Framework (Table 6.2)
%%   6.4 Test Scenarios (Table 6.3)
%%   6.5 Metrics and KPIs (Tables 6.4 and 6.5)
%%   6.6 Testing Infrastructure (Table 6.6)
%%   6.7 Validity and Limitations
%%   6.8 Summary
%%
%% External labels referenced:
%%   chap:research-questions (Ch 3), chap:results (Ch 7),
%%   chap:discussion (Ch 8),
%%   subsec:evaluation-philosophy (4.1.4),
%%   tab:challenges (Table 5.3), app:security (Appendix D),
%%   subsec:future-work-external-validation (9.3.1 in Table 9.1)
%%
%% Citations: none (methodology chapter, not literature-based).
%% ======================================================================

\chapter{Evaluation and Testing}\label{chap:evaluation}

This chapter describes the evaluation methodology employed to assess CCRMS against the research questions outlined in Chapter~\ref{chap:research-questions}. Emphasis is placed on \textbf{how} the system was evaluated; the results are presented in Chapter~\ref{chap:results}. This chapter encompasses the evaluation approach, the four-layer testing architecture, the eight test scenarios, the 11 KPIs, the testing infrastructure, and the validity considerations that inform our interpretation of the findings. The measurement-focused philosophy articulated in Section~\ref{subsec:evaluation-philosophy} informs our methodology: KPI targets serve as reference points rather than pass/fail criteria.

\section{Evaluation Approach}\label{sec:eval-approach}

Our evaluation methodology incorporates unit testing, integration testing, performance benchmarking, security assessment, comparative analysis against a centralized baseline, and Monte Carlo simulation. This comprehensive approach is essential because the research questions span the technical, economic, and security domains, and no single method can address all three. We triangulate quantitative metrics with qualitative insights derived from literature synthesis. Table~\ref{tab:eval-mapping} maps each evaluation activity in relation to the specific research sub-questions it addresses.

\begin{longtable}{@{}L{0.28\linewidth} L{0.36\linewidth} L{0.30\linewidth}@{}}
\caption{Mapping of Evaluation Activities to Research Sub-Questions}\label{tab:eval-mapping}\\
\toprule
\textbf{Research Question} & \textbf{Primary evaluation activity} & \textbf{Supporting activities} \\
\midrule
\endfirsthead

\multicolumn{3}{l}{\emph{Table~\ref{tab:eval-mapping} continued from previous page}}\\
\toprule
\textbf{Research Question} & \textbf{Primary evaluation activity} & \textbf{Supporting activities} \\
\midrule
\endhead

\midrule
\multicolumn{3}{r}{\emph{Continued on next page}}\\
\endfoot

\bottomrule
\endlastfoot

SQ1 (XCM cross-chain rights operations and finality) & XCM latency benchmark, end-to-end XCM test & Unit tests for \texttt{xcm\_*} extrinsics \\
\addlinespace
SQ2 (Latency vs decentralization trade-offs) & Throughput and latency benchmark, HHI calculation & Comparative analysis vs centralized \\
\addlinespace
SQ3 (Unified rights token with royalty distribution) & Unit tests, royalty propagation tests, comparative benchmark & Monte Carlo creator revenue simulation \\
\addlinespace
SQ4 (Security and authorization) & Security audit, authorization unit tests & Reliability test \\
\addlinespace
SQ5 (Economic viability) & Monte Carlo creator revenue simulation & Literature comparison \\
\end{longtable}

Each KPI (Section~\ref{sec:kpis}) has a specified target derived from the Phase~2 methodology; however, these targets are regarded as reference points rather than definitive pass/fail criteria. All 11 KPIs met their targets; the methodology's validity would have been maintained even if some targets had not been met.

\section{Testing Layers}\label{sec:testing-layers}

We organize testing into four layers, each executed at the lowest applicable level of abstraction. Figure~\ref{fig:testing-layers} summarises the stack; the subsections below describe each layer.

\begin{figure}[htbp]
\vspace{\baselineskip}
\centering
\begin{tikzpicture}[
    layer/.style={draw, rectangle, minimum width=12cm, text width=11.5cm,
                  minimum height=1.4cm, align=center, font=\small},
    node distance=3pt
]
\node[layer, fill=green!18] (l1) at (0,0) {
    \textbf{Layer 1 --- Unit Tests}\\
    {\footnotesize 56 tests in mock runtime: 23 functional · 10 security · 11 XCM · 5 royalty · 5 auto-renewal · 2 metadata;
     \texttt{cargo test}; $\sim$3\,s; run on every commit.
     \emph{Does not cover:} network latency, live cross-chain routing, or sustained load.}
};
\node[layer, fill=blue!18, above=3pt of l1] (l2) {
    \textbf{Layer 2 --- XCM Simulator}\\
    {\footnotesize In-process two-parachain dispatch via \texttt{xcm-simulator}; no Zombienet required.
     Isolates XCM message construction and configuration failures from block-timing and network conditions.
     \emph{Does not cover:} wall-clock latency or live relay-chain routing.}
};
\node[layer, fill=orange!18, above=3pt of l2] (l3) {
    \textbf{Layer 3 --- Zombienet End-to-End}\\
    {\footnotesize Live Rococo-local relay chain; Para~A and Para~B connected via HRMP;
     7 automated \texttt{@polkadot/api} benchmark scripts producing JSON output.
     Primary evidence source for KPIs 1--8.
     \emph{Does not cover:} real-network latency or multi-collator topologies.}
};
\node[layer, fill=purple!18, above=3pt of l3] (l4) {
    \textbf{Layer 4 --- Snowbridge End-to-End}\\
    {\footnotesize Full bridge stack: 4 parachains, Geth v1.17.1, Lodestar v1.35.0, 16 Gateway contracts,
     beacon and execution relayers; $\approx$40\,min beacon synchronisation.
     Validates the complete Ethereum-to-Polkadot bridge pathway; run infrequently.}
};
\end{tikzpicture}
\caption{Four-layer testing architecture (Layer~1 at base: fastest and most granular, run on every commit; Layer~4 at apex: broadest scope, run infrequently). Each layer tests what lower layers cannot: unit tests isolate individual extrinsic logic; the XCM simulator isolates message-construction failures; Zombienet provides definitive wall-clock performance evidence; Snowbridge validates the complete Ethereum bridge pathway.}
\label{fig:testing-layers}
\end{figure}

\textbf{Layer 1: Unit tests} (56 tests in a mock runtime). The suite encompasses 23 functional tests, ten security tests, 11 XCM tests, five royalty propagation verifications, five auto-renewal tests, and two metadata query tests. We run them with \texttt{cargo test}; all 56 tests pass on every commit in approximately three seconds. The unit layer provides the most accurate indications of individual extrinsic correctness but does not cover cross-chain interactions, network latency, or sustained load testing.

\textbf{Layer 2: XCM simulator tests.} The \texttt{xcm-simulator} framework dispatches XCM messages between two parachains in-process, without the need to deploy a Zombienet network. This design isolates cross-chain logic from operational complexity: a simulator failure indicates a message construction or XCM configuration, whereas a Zombienet failure may be due to network conditions or block timing.

\textbf{Layer 3: Zombienet end-to-end tests.} Live multi-chain tests on a Rococo-local relay chain with two validators (Alice, Bob) and two CCRMS parachains (Para~A, Para~B) connected via HRMP. JavaScript scripts in \texttt{scripts/perf/} submit extrinsics via \texttt{@polkadot/api}, measure wall-clock time from submission to block inclusion, and record results as JSON. This layer provides the definitive evidence for cross-chain latency and exercises sustained-load and stress scenarios.

\textbf{Layer 4: Snowbridge end-to-end tests.} Extends the Zombienet topology with the full Snowbridge stack, comprising four parachains, Geth, Lodestar, 16 Solidity Gateway contracts, and two relayer processes. Tests confirm the functionality of the complete Ethereum-to-Polkadot bridge pathway. Two 1~ETH transactions were successfully bridged during Phase~4. This layer necessitates approximately 40 minutes of beacon synchronization and is engaged infrequently, primarily for comprehensive bridge validation.

\section{Performance Testing Framework}\label{sec:perf-framework}

Seven automated scripts in \texttt{scripts/perf/} cover various aspects of system behavior, each generating machine-readable JSON output. Table~\ref{tab:benchmarks} summarizes the benchmark categories.

\begin{longtable}{@{}L{0.16\linewidth} L{0.28\linewidth} L{0.50\linewidth}@{}}
\caption{Benchmark Scripts}\label{tab:benchmarks}\\
\toprule
\textbf{Category} & \textbf{Script} & \textbf{What it measures} \\
\midrule
\endfirsthead

\multicolumn{3}{l}{\emph{Table~\ref{tab:benchmarks} continued from previous page}}\\
\toprule
\textbf{Category} & \textbf{Script} & \textbf{What it measures} \\
\midrule
\endhead

\midrule
\multicolumn{3}{r}{\emph{Continued on next page}}\\
\endfoot

\bottomrule
\endlastfoot

Throughput & \texttt{local-throughput.mjs} & TPS across 11 extrinsic types at batch sizes 1, 10, 20, 50, 100 \\
\addlinespace
Stress & \texttt{stress-test.mjs} & Sustained TPS from 200 concurrent accounts over 180 seconds \\
\addlinespace
XCM latency & \texttt{xcm-latency.mjs} & Block delta between send (Para~B) and event arrival (Para~A) \\
\addlinespace
Block weight & \texttt{block-utilization.mjs} & Percentage of \texttt{ref\_time} and \texttt{proof\_size} consumed at batch sizes 50--300 \\
\addlinespace
Storage & \texttt{storage-growth.mjs} & SCALE-encoded byte size per content item, subscriber, and view pack \\
\addlinespace
Resource monitoring & \texttt{resource-monitor.mjs} & CPU, memory, and disk usage via Prometheus scraping during load \\
\addlinespace
Reliability & \texttt{reliability-test.mjs} & Baseline uptime, MTTR after SIGTERM, state persistence, recovery uptime \\
\addlinespace
Comparative & \texttt{centralized-drm-benchmark/} & Same six operations on Express.js + SQL.js under identical conditions \\
\end{longtable}

A separate Monte Carlo simulation (\texttt{monte-carlo-revenue.mjs}) models 1,000 creators across 10,000 iterations under four platform configurations (centralized 30--45\% fees, Web3 marketplace 5--10\%, bridge-based 8--15\%, CCRMS 1--5\%), incorporating algorithmic discovery, organic growth, and monthly churn. This provides the primary evidence for research question SQ5. We document the methodology in \texttt{scripts/perf/results/monte-carlo-findings.md}.

All scripts are reproducible: each is an independent Node.js program with explicit parameters and deterministic output paths, archived as timestamped JSON in \texttt{scripts/perf/results/}. Zombienet topologies and SDK versions are version-controlled.

\section{Test Scenarios}\label{sec:test-scenarios}

We designed eight test scenarios to cover the full range of evaluation activities. Table~\ref{tab:scenarios} summarizes each scenario; the corresponding scripts live in \texttt{scripts/perf/}.

\begin{longtable}{@{}L{0.03\linewidth} L{0.17\linewidth} L{0.45\linewidth} L{0.28\linewidth}@{}}
\caption{Test Scenarios}\label{tab:scenarios}\\
\toprule
\textbf{\#} & \textbf{Scenario} & \textbf{Description} & \textbf{Target} \\
\midrule
\endfirsthead

\multicolumn{4}{l}{\emph{Table~\ref{tab:scenarios} continued from previous page}}\\
\toprule
\textbf{\#} & \textbf{Scenario} & \textbf{Description} & \textbf{Target} \\
\midrule
\endhead

\midrule
\multicolumn{4}{r}{\emph{Continued on next page}}\\
\endfoot

\bottomrule
\endlastfoot

1 & Throughput and latency & Submit concurrent batches (1, 10, 20, 50, 100) for each of the 11 local extrinsics; measure submission-to-inclusion time. & $>$10 TPS sustained; $\sim$6~s inclusion latency \\
\addlinespace
2 & Block weight saturation & Submit progressively larger \texttt{register\_content} batches (50--300); query \texttt{system.blockWeight} at inclusion. & Block weight $<$75\% of budget at max batch \\
\addlinespace
3 & Storage growth & Register content, subscribers, and view packs in incremental quantities; measure SCALE-encoded byte size per map. & $<$500 bytes per content item; linear growth \\
\addlinespace
4 & XCM latency & Send \texttt{xcm\_subscribe}, \texttt{xcm\_purchase\_views}, \texttt{xcm\_purchase\_ownership} from Para~B to Para~A via the canonical \texttt{WithdrawAsset} $\to$ \texttt{BuyExecution} $\to$ \texttt{Transact} sequence; measure block delta. & $<$10 parachain blocks per operation \\
\addlinespace
5 & Stress test & 200 funded accounts continuously submit \texttt{register\_content} for 180 seconds. & Sustained TPS substantially above single-batch throughput \\
\addlinespace
6 & Security audit & Static analysis (\texttt{cargo clippy}, \texttt{cargo audit}), manual review of all 17 extrinsics, targeted unit tests for authorization gaps and edge cases. & 0 Critical / 0 High vulnerabilities; complete authorization coverage \\
\addlinespace
7 & Reliability & Monitor block production for 120~s (baseline), SIGTERM the collator, restart, verify state, monitor for 120~s (recovery). & MTTR $<$60~s; 100\% state persistence \\
\addlinespace
8 & Comparative analysis & Re-implement the same six operations in Express.js + SQL.js (in-memory SQLite); run identical benchmarks against both. & Document the performance cost of decentralization \\
\end{longtable}

The comparative benchmark is intentionally minimal, using Express.js for HTTP and SQL.js for in-process SQLite. A more comprehensive, centralized implementation incorporating caching, replication, and distributed databases would exhibit lower performance than our baseline but would remain significantly faster than CCRMS. The minimal version serves as a near-optimal reference point for comparison and precludes any implication of manipulation.

\section{Metrics and Key Performance Indicators}\label{sec:kpis}

Our evaluation employs 11 key performance indicators (KPIs) across four domains: technical, economic, security, and reliability. Each KPI is defined by its description, measurement methodology, originating source (the test scenario or activity responsible for generating it), and a target value. These targets are derived from the Phase~2 evaluation methodology and are regarded as reference points rather than strict pass/fail criteria.

\begin{longtable}{@{}L{0.03\linewidth} L{0.22\linewidth} L{0.36\linewidth} L{0.17\linewidth} L{0.16\linewidth}@{}}
\caption{KPI Definitions}\label{tab:kpi-defs}\\
\toprule
\textbf{\#} & \textbf{KPI} & \textbf{Definition} & \textbf{Target} & \textbf{Source} \\
\midrule
\endfirsthead

\multicolumn{5}{l}{\emph{Table~\ref{tab:kpi-defs} continued from previous page}}\\
\toprule
\textbf{\#} & \textbf{KPI} & \textbf{Definition} & \textbf{Target} & \textbf{Source} \\
\midrule
\endhead

\midrule
\multicolumn{5}{r}{\emph{Continued on next page}}\\
\endfoot

\bottomrule
\endlastfoot

1 & \textbf{Throughput (TPS)} & Successful transactions per second under sustained load & $>$10 TPS & Stress test \\
\addlinespace
2 & \textbf{Inclusion latency} & Time from submission to block inclusion & $<$12 seconds (2 blocks) & Throughput test \\
\addlinespace
3 & \textbf{XCM latency} & Block delta for cross-chain operations & $<$10 blocks & XCM latency test \\
\addlinespace
4 & \textbf{Block utilization} & Percentage of block weight consumed at capacity & $<$75\% & Block saturation test \\
\addlinespace
5 & \textbf{Storage efficiency} & Bytes per content item and per subscriber & $<$500 bytes & Storage growth test \\
\addlinespace
6 & \textbf{Uptime} & Percentage of expected blocks actually produced & $>$85\% & Reliability test \\
\addlinespace
7 & \textbf{MTTR} & Time from crash to resumed block production & $<$60 seconds & Reliability test \\
\addlinespace
8 & \textbf{State persistence} & Data integrity after crash restart & 100\% & Reliability test \\
\addlinespace
9 & \textbf{Security findings} & Critical/High vulnerabilities in dissertation code & 0 Critical, 0 High & Security audit \\
\addlinespace
10 & \textbf{Test coverage} & Error variants exercised by unit tests & $>$80\% & Test suite \\
\addlinespace
11 & \textbf{Decentralization ratio} & Performance cost vs centralized baseline & Documented & Comparative analysis \\
\end{longtable}

\begin{longtable}{@{}L{0.03\linewidth} L{0.22\linewidth} L{0.18\linewidth} L{0.33\linewidth} L{0.18\linewidth}@{}}
\caption{KPI Results Summary}\label{tab:kpi-results}\\
\toprule
\textbf{\#} & \textbf{KPI} & \textbf{Target} & \textbf{Measured} & \textbf{Status} \\
\midrule
\endfirsthead

\multicolumn{5}{l}{\emph{Table~\ref{tab:kpi-results} continued from previous page}}\\
\toprule
\textbf{\#} & \textbf{KPI} & \textbf{Target} & \textbf{Measured} & \textbf{Status} \\
\midrule
\endhead

\midrule
\multicolumn{5}{r}{\emph{Continued on next page}}\\
\endfoot

\bottomrule
\endlastfoot

1 & Throughput & $>$10 TPS & \textbf{26.2 TPS} sustained & \textbf{Pass} \\
\addlinespace
2 & Inclusion latency & $<$12~s & \textbf{$\sim$6~s} (1 block) & \textbf{Pass} \\
\addlinespace
3 & XCM latency & $<$10 blocks & \textbf{3--5 blocks} & \textbf{Pass} \\
\addlinespace
4 & Block utilization & $<$75\% & \textbf{12.9\%} at 300 txs & \textbf{Pass} \\
\addlinespace
5 & Storage efficiency & $<$500 bytes & \textbf{191 bytes}/content, \textbf{112 bytes}/subscriber & \textbf{Pass} \\
\addlinespace
6 & Uptime & $>$85\% & \textbf{90\%} & \textbf{Pass} \\
\addlinespace
7 & MTTR & $<$60~s & \textbf{$\sim$15--20~s} & \textbf{Pass} \\
\addlinespace
8 & State persistence & 100\% & \textbf{100\%} & \textbf{Pass} \\
\addlinespace
9 & Security findings & 0 Critical/High & \textbf{0 Critical, 0 High, 0 unresolved Medium} & \textbf{Pass} \\
\addlinespace
10 & Test coverage & $>$80\% & \textbf{84\%} (16/19 error variants) & \textbf{Pass} \\
\addlinespace
11 & Decentralization ratio & Documented & \textbf{$270\times$ slower, qualitatively superior} & \textbf{Documented} \\
\end{longtable}

All 11 KPIs achieved their respective targets. The security audit identified eight findings (see Table~\ref{tab:challenges}, Appendix~\ref{app:security}): two vulnerabilities were remediated (B, E), three design properties were accepted (A, C, D), and three observations relating to positive reliability were noted (F, G, H). Chapter~\ref{chap:discussion} offers an interpretation of these results.

\section{Testing Infrastructure}\label{sec:test-infra-ch6}

\begin{longtable}{@{}L{0.28\linewidth} L{0.65\linewidth}@{}}
\caption{Test Environment Specifications}\label{tab:test-env}\\
\toprule
\textbf{Component} & \textbf{Specification} \\
\midrule
\endfirsthead

\multicolumn{2}{l}{\emph{Table~\ref{tab:test-env} continued from previous page}}\\
\toprule
\textbf{Component} & \textbf{Specification} \\
\midrule
\endhead

\midrule
\multicolumn{2}{r}{\emph{Continued on next page}}\\
\endfoot

\bottomrule
\endlastfoot

Machine & macOS Darwin 25.3.0, Apple Silicon \\
Relay chain & Rococo-local, 2 validators (alice, bob) \\
Parachain A & Content Rights (parachain 100), 1 collator \\
Parachain B & Consumer chain (parachain 200), 1 collator \\
Bridge Hub & Parachain 1013 (Snowbridge tests only) \\
AssetHub & Parachain 1000 (Snowbridge tests only) \\
Ethereum & Geth v1.17.1 + Lodestar v1.35.0 (Snowbridge tests only) \\
Zombienet provider & Native \\
Block time & $\sim$6 seconds (parachain) \\
Node binary & \texttt{parachain-template-node} (release build) \\
Polkadot SDK version & stable2512 (December 2024) \\
Test scripts & \texttt{scripts/perf/*.mjs} (Node.js, \texttt{@polkadot/api} library) \\
Centralized benchmark & \texttt{centralized-drm-benchmark/} (Express.js + SQL.js) \\
\end{longtable}

We designed the test environment to be minimalistic, comprising a single development machine that operates a local Zombienet network. This methodology simplifies the configuration process and removes dependency on external infrastructure, although it has implications for the interpretation of the results, as discussed in Section~\ref{sec:validity}. The benchmark scripts and raw JSON outputs are maintained under version control within the repository, thereby allowing any reader with access to the Polkadot SDK toolchain to independently re-execute them.

\section{Validity and Limitations}\label{sec:validity}

\textbf{Internal validity.} We conduct all benchmarks multiple times and report the average results; maintain a consistent test environment across all runs; apply an identical methodology to CCRMS and the centralized baseline; and execute the deterministic unit test suite on each commit. Therefore, any variation across runs can be attributed to the system under test rather than environmental fluctuations.

\textbf{External validity.} This represents the dimension in which our evaluation is most limited. Four constraints are applicable: (1) we take measurements on a local Zombienet environment with two relay-chain validators and without real network latency, thus the figures serve as upper bounds within the test environment rather than predictive metrics for production; (2) the Snowbridge integration uses a local Ethereum stack instead of a public testnet, avoiding variability in block times, MEV, contention, and beacon chain reorganizations; (3) because of the maturation gap of Polkadot~2.0, the latency figures reflect a pre-2.0 toolchain, and several measurements are expected to improve significantly in a production Polkadot~2.0 environment; (4) our scale is modest, comprising up to 200 concurrent accounts and a few hundred content items, and empirical verification of linear storage growth at production scale has not yet been conducted.

\textbf{Construct validity.} The figure of 26.2 TPS quantifies the successful inclusion rate under sustained load and constitutes an essential yet insufficient element of user-perceived performance, which also relies on wallet user experience, mempool latency, and RPC responsiveness. The baseline uptime of 90\% assesses single-collator block production and does not account for multi-collator production availability. The Monte Carlo simulation evaluates \textbf{expected creator revenue under a probabilistic model of platform behavior}, rather than actual revenue in operational environments; the validity of this construct is contingent upon the input parameter ranges (derived from publicly available data on YouTube, Spotify, and OpenSea) and the underlying model assumptions (such as algorithmic discovery, churn, and fees), which are explicitly documented in the findings.

\textbf{Threats to validity.} (a)~\textbf{Selection bias:} the eight test scenarios may not encompass production conditions such as high subscriber concurrency on a single content item, large multi-collator topologies, or extended multi-month uptime. (b)~\textbf{Measurement instrument bias:} any bug in \texttt{@polkadot/api} or \texttt{system.blockWeight} would propagate into our results; while the library is extensively used, it is not devoid of bugs. (c)~\textbf{Confirmation bias:} the author was responsible for designing, implementing, and evaluating the system; this risk is mitigated through the use of externally referenced KPI targets and by disclosure of all scripts and raw results. (d)~\textbf{Single-author validation:} no external code review, third-party security audit, or benchmark validation has been conducted; we recognize this as a significant limitation and intend to address it as future work (Section~\ref{subsec:future-work-external-validation}).

\section{Summary}\label{sec:chapter-6-summary}

The evaluation employs four testing layers, eight test scenarios, seven automated benchmarks, a Monte Carlo simulation, and 11 KPIs grounded in a measurement-focused philosophy. Section~\ref{sec:validity} acknowledges the primary validity constraints. Chapter~\ref{chap:results} presents the results.
